\section{Implementation for MobileNetV2}
\label{app:implementation-mobilenetv2}

\subsection{Iterative Pruning}
\label{sec:mbv2-iterative-implementation}
The iterative run uses Algorithm~\ref{alg:iterative-pruning} with
the phase budgets in Table~\ref{tab:pruning-schedules}. Search and recovery restart from ImageNet
weights; only the removed-channel set passes between phases. Each recovery
phase contributes one validation-selected model to the trajectory.

We use MobileNetV2~\cite{sandler2018mobilenetv2} with its standard
channel widths (2{,}480{,}072 parameters), using torchvision
\texttt{IMAGENET1K\_V2} weights and a new 200-way classifier.
Learning rates are $10^{-4}$ for ImageNet-initialized weights and
$10^{-3}$ for gates and classifier. The dense Baseline is fine-tuned
for 45 epochs using the same optimizer settings.
% Source: mobilenetv2_tinyimagenet200_224_pretrained.yaml; adamw_finetune.yaml.

\paragraph{Gate placement.}
We gate the hidden width of all 16 inverted residual blocks with
an expansion layer, after depthwise convolution, BatchNorm and
ReLU6 (ReLU clipped at 6). Each hidden gate couples one expansion filter, one depthwise
filter, their BatchNorm parameters and the corresponding projection
input channel. Placing the gate after normalization makes masking
equivalent to physical removal; placing it before BatchNorm would
allow a closed channel to become nonzero. The ten blocks with identity
shortcuts also have output gates after projection BatchNorm:
\begin{equation}
\mathbf{x}_l=\mathbf{x}_{l-1}
+\mathbf{g}^{\mathrm{out}}_l\odot
b_{\theta,l}(\mathbf{g}_l\odot a_{\theta,l}(\mathbf{x}_{l-1})).
\end{equation}
Closing an output gate removes the corresponding projection filter
and BatchNorm parameters while the shortcut preserves that channel
and the block output width. Hidden and output gates follow the same
objective and pruning rule: 7{,}104 hidden gates and 792 output gates,
7{,}896 in total.

Each gate has an input-independent learnable logit with small random
initialization perturbations. Training uses sampled ST masks;
evaluation thresholds the opening probabilities. Initialization and
evaluation settings are listed in Table~\ref{tab:method-hyperparameters}.
The sampling temperature is $T=1$. Thus the calibration average in the Methodology
section is simply $\overline r_{l,c}=r_{l,c}$.
% Sources: MaskedGumbelLayer; MaskedGumbelInvertedResidual;
% PrunedGumbelInvertedResidual; ours_iterative_channels configuration.

\paragraph{Objective and normalization.}
Let $M$ be the number of gate groups and $\mathcal C_j$ the
$C_j>0$ channels in group $j$ not
permanently masked in the current search phase. Hidden and output
gates form separate groups. We average within groups, then across them:
\begin{equation}
S(r_\phi)=\frac{1}{M}\sum_{j=1}^{M}\frac{1}{C_j}
\sum_{c\in\mathcal C_j}r_{j,c}.
\label{eq:mbv2-regularizer-means}
\end{equation}
The implemented objective is
\begin{equation}
\mathcal J_{\mathrm{impl}}
=\mathcal L_{\mathrm{CE}}+\lambda S(r_\phi).
\label{eq:mbv2-implemented-objective}
\end{equation}
Each group has equal weight, giving coefficient $\lambda/(MC_j)$
to each channel. The reported coefficients use this normalization.

\paragraph{Adaptive sparsity.}
The controller holds $\lambda$ fixed during warmup, then updates it at
the interval and within the bounds in Table~\ref{tab:method-hyperparameters}.
The reference is a separate dense model trained for 30 epochs, distinct
from the 45-epoch Baseline. Its accuracy at epoch $t$ is compared with
the moving average of gated accuracy. The window and accuracy tolerances
are specified in Table~\ref{tab:method-hyperparameters}. We increase
$\lambda$ when the gap is at most $\delta_s$, decrease it above
$\delta_h$, and hold it otherwise.
The log-space step is
$\eta=\log(1.25)\,2^k$, $k\in\{0,1,2\}$. During increases, $k$
rises by one (up to 2) if mean closing probability $1-r_{l,c}$ over
all gates grows by less than 0.02 since the previous update, resets for growth in
$[0.02,0.07]$ and otherwise stays unchanged; it also resets whenever
$\lambda$ is not increased.
% Source: adaptive_lambda configuration; training/adaptive_lambda.py.

\paragraph{Channel selection.}
Algorithm~\ref{alg:iterative-pruning} uses the threshold and cap in
Table~\ref{tab:method-hyperparameters}.
Because gates start below the pruning threshold, search must raise their
scores to protect them, subject to the parameter cap. Recovery evaluates
the physically compact network without gates.

\subsection{One-Pass Pruning}
\label{sec:mbv2-one-shot-implementation}
The seven settings use the same training, gate, controller and
initialization protocol as iterative MobileNetV2, with one search
and one recovery phase. They vary only the selection threshold and cap;
the seven pairs are listed in Table~\ref{tab:mbv2-tin-all}, with phase
budgets in Table~\ref{tab:pruning-schedules}.
This is one cycle of Algorithm~\ref{alg:iterative-pruning}; it does not
use the ResNet checkpoint selector in Algorithm~\ref{alg:accuracy-guided-pruning}.
% Phase budgets and protocol equivalence confirmed by the author, 2026-11-11.
